%!TEX root = ../../../main.tex

\chapter{Datenbasis und Vorverarbeitung}
\label{chap:02-data-preprocessing}

% Bewertungskriterium: Vorverarbeitung & Datenstrukturen (15 %)

Die 30 PDFs mit zusammen 34 Seiten haben alle eine Textebene, sind aber gestaltet wie echte Post. Die Tabellen haben keine Rahmenlinien, nur abwechselnd hinterlegte Zeilen, und fast jede Seite endet mit Bank- und Registerangaben. Fünf Diagramme enthalten Werte, die nur im Bild stehen, etwa den Stromverbrauch 2022 mit 3.480\,kWh.

Der \texttt{DocumentParser} (PyMuPDF) sucht auf jeder Seite zuerst die Tabellen und nimmt ihre Fläche aus dem Fließtext heraus, damit kein Wert doppelt vorkommt. Wie unzuverlässig das anfangs war, zeigte erst der Vergleich mit den PDFs. Je nach Strategie verlor \texttt{find\_tables()} Kopfzeilen oder fasste in Dokument~29 Rechnungsbetrag und Abschläge zu einer Zeile zusammen. Der Parser lässt deshalb zwei Strategien laufen und behält die Tabelle mit mehr Zeilen. Den Zelltext von PyMuPDF übernimmt er nicht, weil Zahlen darin zerrissen sind (\enquote{1~189~00~EUR}), sondern setzt die Zellen selbst aus den Textspans zusammen. Eine fette erste Zeile gilt als Spaltenkopf, und jede weitere Zeile wird mit den Spaltennamen als Text gespeichert. Überschriften erkennt der Parser an Fettdruck und Schriftgröße und legt sie im Überschriften-Stapel (\texttt{heading\_stack}) ab, Kopf- und Fußzeilen verwirft er. Verworfen wurden feste Chunk-Längen, weil sie Tabellen zerschnitten hätten, und ein Layoutmodell wie Docling. Es bringt eigene KI-Modelle mit, und warum es eine Stelle als Tabelle erkennt, ließe sich nicht Regel für Regel erklären.

Ein Durchlauf über alle PDFs dauert etwa 31\,s und ergibt 564 Chunks: 178 \texttt{text}, 82 \texttt{heading}, 299 \texttt{table\_row} aus 65 Tabellen und 5 \texttt{image}. Listing~\ref{lst:chunk-json} zeigt das JSON-Zwischenformat an der Tabellenzeile mit dem falschen Arbeitspreis aus Kapitel~\ref{chap:01-problem-use-case}. Die Metadaten stammen nicht aus dem PDF, sondern aus der von Hand gepflegten Datei \texttt{documents.yaml}, und die Suche filtert später danach (Kapitel~\ref{chap:04-retrieval}).

\begin{lstlisting}[language={},numbers=none,columns=fullflexible,basicstyle=\ttfamily\scriptsize,aboveskip=8pt,caption={Tabellenzeile der Stromrechnung 2025 im Zwischenformat (gekürzt)},label={lst:chunk-json},literate={ö}{{\"o}}1 {ü}{{\"u}}1 {Ö}{{\"O}}1]
{"chunk_id": "06_Strom_Jahresabrechnung_2025_SuedVolt_p1_tab03_r01",
 "chunk_type": "table_row", "page_number": 1,
 "content": "Abrechnungsposition: Arbeitspreis ÖkoStrom | ... | Preis Brutto: 39,50 ct/kWh | ...",
 "heading_stack": ["Ihre Jahresrechnung Strom für den Zeitraum 01.01.2025 bis 31.12.2025",
                   "2. Detailberechnung der Energiekosten (Tarif: SüdVolt ÖkoStrom Klassik 12)"],
 "metadata": {"utility_type": "electricity", "doc_type": "bill", "contract_id": "SV-2024-99104",
              "billing_period_start": "2025-01-01", "billing_period_end": "2025-12-31", ...}}
\end{lstlisting}

Im Enrichment schreibt \texttt{gemma3:4b} zu jedem Diagramm eine Bild-Caption und \texttt{llama3.2} zu jeder der 65 Tabellen eine Tabellen-Summary als eigenen Chunk. Damit sind es 629 Chunks. Weil ein Lauf gut 26 Minuten dauert, liegen die Modellantworten in einem Cache, dessen Schlüssel ein SHA-256-Hash über Modell, Prompt, Optionen und gegebenenfalls die Bildbytes ist. Im ersten Lauf waren 21 der 65 Summaries am Token-Limit abgeschnitten, und zu einem Diagramm ohne beschriftete Säulen erfand das Vision-Modell Monatswerte. Mit angepassten Prompts verschwand beides, inhaltliche Fehler blieben aber. So vergleicht die Summary zur Gebührenanpassung Wasser 2026 den alten Brutto- mit dem neuen Nettopreis und meldet eine Senkung, obwohl der Bruttopreis steigt. Summaries dienen deshalb nur der Suche, zitiert und gerechnet wird mit den Tabellenzeilen. Ohne sie fand die Suche bei 8 statt 9 der 10 Retrieval-Testfragen (Kapitel~\ref{chap:04-retrieval}) das erwartete Dokument.

Die Regeln sind auf dieses eine Layout abgestimmt, und Briefköpfe hält der Parser teilweise für Tabellen oder Überschriften. In der Fehleranalyse (Kapitel~\ref{chap:05-evaluation}) ging trotzdem keine falsche Antwort auf einen fehlenden oder entstellten Chunk zurück.
